% !TeX program = xelatex
% ============================================================================
%  第 6 课　复合作为运算：S₃ 是个群　—— 学生版讲义
%  与 02-课件/第06课-复合作为运算/第06课-复合作为运算.tex 同步
%  记号：复合 xy 读作「先做 x，再做 y」
% ============================================================================

\zhipianye{6}

\xhlesson{6}{复合作为运算：\texorpdfstring{\(S_3\)}{S3} 是个群}{}

\begin{mubiao}
  \begin{itemize}
    \item 能把「先做一个再做另一个」当成一种运算（复合）。
    \item 会一条一条验证这 \(6\) 个动作满足群的四条规则。
    \item 能看懂乘法表，并发现「每一行都是一次重排」。
  \end{itemize}
\end{mubiao}

\section*{一、给六个动作起名字}

\begin{example}
  \begin{center}
  \begin{tabular}{@{}clc@{}}
    \toprule
    名字 & 动作 & 两行式 \\
    \midrule
    \(e\) & 不动 & \(\begin{pmatrix}A&B&C\\A&B&C\end{pmatrix}\) \\[6pt]
    \(r\) & 顺时针转 \(120^\circ\) & \(\begin{pmatrix}A&B&C\\B&C&A\end{pmatrix}\) \\[6pt]
    \(r^2\) & 顺时针转 \(240^\circ\) & \(\begin{pmatrix}A&B&C\\C&A&B\end{pmatrix}\) \\[6pt]
    \(f\) & 沿过 \(A\) 的轴翻 & \(\begin{pmatrix}A&B&C\\A&C&B\end{pmatrix}\) \\
    \bottomrule
  \end{tabular}
  \end{center}
\end{example}

\begin{definition}[复合]
  两个动作 \(x\) 和 \(y\)，先做 \(x\)、再做 \(y\)，合起来还是一个动作，记作 \(xy\)。
  这个动作叫做 \(x\) 与 \(y\) 的\textbf{复合}。
\end{definition}

\begin{definition}[为什么也管它叫「乘法」]
  你看 \(xy\) 写成两个字母挨在一起，跟\textbf{乘法省掉乘号}的样子一模一样——
  第 2 课讲过：\(3a\) 就是 \(3\times a\)，乘号可以省掉。\\[3pt]
  所以数学家也常把复合叫做「\textbf{乘法}」，把 \(xy\) 读作「\(x\) 乘 \(y\)」。\\[3pt]
  但要留个心眼：这个「乘法」\textbf{一般不能交换}——就是第 4 课见过的那件事，
  下面 \(rf\) 和 \(fr\) 算出来的结果就不一样。
\end{definition}

\section*{二、先做哪个，结果一样吗}

\begin{example}[先看一个示范]
  算 \(rf\)，就是「先做 \(r\)，再做 \(f\)」，一路跟着 \(A\) 走：
  \[
    A \xrightarrow{\ r\ } B \xrightarrow{\ f\ } C .
  \]
  \(A\) 先被 \(r\) 送到 \(B\)，再被 \(f\)（沿过 \(A\) 的轴翻）送到 \(C\)，所以 \(A\to B\to C\)。
  剩下的 \(B\)、\(C\) 照着同样的走法自己推。
\end{example}

\begin{dongshou}[跟老师一起算]
  用两行式一路推下去（先 \(x\)，后 \(y\)，就是接着往下走）：

  \(rf\)：\(A\to\underline{\hspace{1.6em}}\to\underline{\hspace{1.6em}}\)，
  \(B\to\underline{\hspace{1.6em}}\to\underline{\hspace{1.6em}}\)，
  \(C\to\underline{\hspace{1.6em}}\to\underline{\hspace{1.6em}}\)。

  \(fr\)：\(A\to\underline{\hspace{1.6em}}\to\underline{\hspace{1.6em}}\)，
  \(B\to\underline{\hspace{1.6em}}\to\underline{\hspace{1.6em}}\)，
  \(C\to\underline{\hspace{1.6em}}\to\underline{\hspace{1.6em}}\)。
  \liukong{5}
  我的结论：\(rf\) 与 \(fr\) 一样吗？\underline{\hspace{3em}}
\end{dongshou}

\begin{example}
  \[
    rf=\begin{pmatrix}A&B&C\\C&B&A\end{pmatrix}\ (\text{沿过 }B\text{ 的轴翻}), \qquad
    fr=\begin{pmatrix}A&B&C\\B&A&C\end{pmatrix}\ (\text{沿过 }C\text{ 的轴翻}).
  \]
  两个结果不一样，所以 \(rf\neq fr\)。
\end{example}

\section*{三、四条规则，一条一条验}

\begin{dongshou}[在每一条后面写上「成立」或「不成立」]

  \begin{enumerate}
    \item \textbf{封闭}：两个动作连着做，还是这 \(6\) 个之一吗？\underline{\hspace{3em}}
    \item \textbf{结合}：\((xy)z\) 和 \(x(yz)\) 都是「先 \(x\)、再 \(y\)、再 \(z\)」吗？\underline{\hspace{3em}}\\
      白话：复合本来就是「接着做」，所以三个动作连着做，先捆哪两个都一样——
      \textbf{结合律天生成立，不用验算}。第 4 课那个 \((a+b)/2\) 才是真的不结合。
    \item \textbf{单位元}：谁是那个「做了等于没做」的动作？\underline{\hspace{3em}}
    \item \textbf{逆元}：填一填——
      \(r\) 的逆是 \underline{\hspace{1.6em}}，\(r^2\) 的逆是 \underline{\hspace{1.6em}}，
      \(f\) 的逆是 \underline{\hspace{1.6em}}。
  \end{enumerate}
\end{dongshou}

\begin{dingli}[三角形动作群]
  正三角形的 \(6\) 个对称动作，连同复合这个运算，构成一个群，记作 \(S_3\)。
  它\textbf{不满足交换律}。
\end{dingli}

\section*{四、填乘法表}

\begin{dongshou}[只填 \(r\) 行和 \(f\) 行]
  格子里的动作是「\textbf{先横排的，再做竖排的}」。

  \begin{center}
  \begin{tabular}{@{}l|cccccc@{}}
    \toprule
     & \(e\) & \(r\) & \(r^2\) & \(f\) & \(rf\) & \(r^2f\) \\
    \midrule
    \rule{0pt}{2.8ex}先 \(r\) & & & & & & \\
    \midrule
    \rule{0pt}{2.8ex}先 \(f\) & & & & & & \\
    \midrule
    \rule{0pt}{2.8ex}先 \(r^2\)（挑战） & & & & & & \\
    \bottomrule
  \end{tabular}
  \end{center}
  前两行是必做；\textbf{第三行「先 \(r^2\)」是挑战行}，做得动就填，做不动就空着。
  （提示：先 \(r^2\) 就是「先转 \(240^\circ\)」。\(r^2\) 再接着做一次 \(r\)，等于转 \(360^\circ\)，就回到原样了。）
  一行一格地写，别串行——先横排的动作，再做竖排的动作。
\end{dongshou}

\begin{sikao}
  把 \(r\) 那一行读一遍：有没有哪个动作出现了两次？有没有哪个动作没出现？
  \liukong{2}
  所以每一行都是\underline{\hspace{8em}}。
\end{sikao}

\begin{example}
  从 \(f\) 那一行的第二个格子可以读出一个关系：
  \[
    fr=r^2f .
  \]
  白话：「先翻后转，等于先\textbf{倒着}转再翻。」
\end{example}

\begin{xiaojie}
  用自己的话回答：为什么说 \(S_3\)「是一个群，但不是交换的群」？
  \liukong{4}
\end{xiaojie}

\noindent\textbf{下节课}：在 \(S_3\) 里圈出能自成一体的小圈——阶与子群。